\section{研究难点及预设解决方案}

\begin{enumerate}
  \item \textbf{确定性边推导的覆盖率与精度平衡。}统一执行轨迹图的安全检查依赖于四类确定性推导器恢复隐式依赖边，若推导规则过于保守则大量关系保持\texttt{UNKNOWN}，增加后续语义补边负担；若推导规则过于激进则可能引入虚假硬边，削弱阻断决策的可靠性。拟在规范化规则设计中严格限定硬边仅来自可唯一重放的确定性证据（精确值匹配、唯一键映射、稳定对象ID版本传播和可寻址状态差分），其余关系统一归入软边或未知边；通过独立的正反测试夹具和字段变异测试持续验证推导器的精度与召回率。

  \item \textbf{任务契约的构造完整性。}任务契约从可信用户请求和工具目录编译生成，若确定性抽取规则未能覆盖用户请求的全部授权意图表达方式，可能导致契约规则集合出现遗漏，使本应可以被授权的动作被错误阻断。拟采用"确定性抽取优先、隔离的一次结构化候选生成为辅"的两阶段构造流程：确定性规则处理可枚举的显式授权表达，仅在存在歧义时触发一次不读取外部观测的隔离候选生成过程，并在实验中通过契约覆盖率和错误阻断率评估构造质量。

  \item \textbf{攻击边定位模型的强监督标签构造。}研究点二的边级分类模型需要SOURCE/SINK/BENIGN标签进行训练和评测，该标签需由benchmark已知注入位置、实际Runtime读取事件和一次性审核的工具影响目录共同生成，不调用LLM且不需逐轨迹人工标注。拟设计严格的标签构造协议和数据就绪门控机制（$DataReady$条件）：轨迹解析、工具目录审核、标签可重复生成、泄漏检查和数据切分五个条件必须同时满足才可进入正式实验；并通过分层人工抽样核验标签精度，以及label shuffle、未来字段和模板ID检查等泄漏审计确保标签可靠性。

  \item \textbf{图-文本联合模型的跨任务泛化。}攻击边定位模型在特定任务族和工具族上训练后，需要在未见过的任务、注入模板和工具族上保持分类与匹配的稳定性。拟在训练数据构建时将工具名映射到工具族和效果类别以减少对具体API名称的记忆，采用场景组隔离的训练/测试划分确保评测的独立性，并通过统一多任务图自监督预训练使图编码器先学习Agent执行结构先验；在实验中按任务族、攻击模板、来源类别、工具族和模型分组分层报告OOD泛化表现，并通过payload-masked诊断和template-OOD检验模型是否依赖攻击模板记忆。
\end{enumerate}
